\documentclass[10pt,a4paper]{amsart}
\usepackage{amsmath,amssymb,amsthm,graphicx}
% Retained minimal preamble for the Stacks Project chapter sites-cohomology.tex (Cohomology on Sites).
% Verbatim excerpt of the chapter preamble: Xy-pic setup (native preamble lines 14-18) and the native
% theorem declarations with the mathematical macros (native preamble lines 154-217). Copied mechanically;
% nothing added, removed or redefined. The wrapper class and the amsmath/amssymb/amsthm/graphicx base
% packages are supplied by the assembler, so the class/externaldocument part of the native preamble is omitted.
% For commutative diagrams we use Xy-pic
\usepackage[all]{xy}

% We use 2cell for 2-commutative diagrams.
\xyoption{2cell}
% Theorem environments.
%
\theoremstyle{plain}
\newtheorem{theorem}[subsection]{Theorem}
\newtheorem{proposition}[subsection]{Proposition}
\newtheorem{lemma}[subsection]{Lemma}

\theoremstyle{definition}
\newtheorem{definition}[subsection]{Definition}
\newtheorem{example}[subsection]{Example}
\newtheorem{exercise}[subsection]{Exercise}
\newtheorem{situation}[subsection]{Situation}

\theoremstyle{remark}
\newtheorem{remark}[subsection]{Remark}
\newtheorem{remarks}[subsection]{Remarks}

\numberwithin{equation}{subsection}

% Macros
%
\def\lim{\mathop{\mathrm{lim}}\nolimits}
\def\colim{\mathop{\mathrm{colim}}\nolimits}
\def\Spec{\mathop{\mathrm{Spec}}}
\def\Hom{\mathop{\mathrm{Hom}}\nolimits}
\def\Ext{\mathop{\mathrm{Ext}}\nolimits}
\def\SheafHom{\mathop{\mathcal{H}\!\mathit{om}}\nolimits}
\def\SheafExt{\mathop{\mathcal{E}\!\mathit{xt}}\nolimits}
\def\Sch{\mathit{Sch}}
\def\Mor{\mathop{\mathrm{Mor}}\nolimits}
\def\Ob{\mathop{\mathrm{Ob}}\nolimits}
\def\Sh{\mathop{\mathit{Sh}}\nolimits}
\def\NL{\mathop{N\!L}\nolimits}
\def\CH{\mathop{\mathrm{CH}}\nolimits}
\def\proetale{{pro\text{-}\acute{e}tale}}
\def\etale{{\acute{e}tale}}
\def\QCoh{\mathit{QCoh}}
\def\Ker{\mathop{\mathrm{Ker}}}
\def\Im{\mathop{\mathrm{Im}}}
\def\Coker{\mathop{\mathrm{Coker}}}
\def\Coim{\mathop{\mathrm{Coim}}}

% Boxtimes
%
\DeclareMathSymbol{\boxtimes}{\mathbin}{AMSa}{"02}

%
% Macros for moduli stacks/spaces
%
\def\QCohstack{\mathcal{QC}\!\mathit{oh}}
\def\Cohstack{\mathcal{C}\!\mathit{oh}}
\def\Spacesstack{\mathcal{S}\!\mathit{paces}}
\def\Quotfunctor{\mathrm{Quot}}
\def\Hilbfunctor{\mathrm{Hilb}}
\def\Curvesstack{\mathcal{C}\!\mathit{urves}}
\def\Polarizedstack{\mathcal{P}\!\mathit{olarized}}
\def\Complexesstack{\mathcal{C}\!\mathit{omplexes}}
% \Pic is the operator that assigns to X its picard group, usage \Pic(X)
% \Picardstack_{X/B} denotes the Picard stack of X over B
% \Picardfunctor_{X/B} denotes the Picard functor of X over B
\def\Pic{\mathop{\mathrm{Pic}}\nolimits}
\def\Picardstack{\mathcal{P}\!\mathit{ic}}
\def\Picardfunctor{\mathrm{Pic}}
\def\Deformationcategory{\mathcal{D}\!\mathit{ef}}
\usepackage{cleveref}
\crefname{originalsection}{original section}{original sections}
\Crefname{originalsection}{Original section}{Original sections}
\makeatletter
\newlabel{section-limp}{{079X}{source}{}{sourcepointer-1}{}}
\newlabel{section-limp@cref}{{[originalsection][0][]079X}{[1][0][]source}{}{}{}}
\makeatother
\crefname{originallemma}{original lemma}{original lemmas}
\Crefname{originallemma}{Original lemma}{Original lemmas}
\makeatletter
\newlabel{sites-lemma-topos-good-site}{{03CI}{source}{}{sourcepointer-2}{}}
\newlabel{sites-lemma-topos-good-site@cref}{{[originallemma][0][]03CI}{[1][0][]source}{}{}{}}
\makeatother
\crefname{originallemma}{original lemma}{original lemmas}
\Crefname{originallemma}{Original lemma}{Original lemmas}
\makeatletter
\newlabel{lemma-square-triangle}{{0EVY}{source}{}{sourcepointer-3}{}}
\newlabel{lemma-square-triangle@cref}{{[originallemma][0][]0EVY}{[1][0][]source}{}{}{}}
\makeatother
\title{A distinguished triangle from a pushout square of presheaves of sets}
\author{The Stacks Project Authors}
\date{stacks-project:sites-cohomology.tex:local-snapshot-20261011; GNU FDL 1.2 or later}
\begin{document}
\maketitle
\noindent\textit{Editorial note.} All delivered mathematical text is the author's own, copied from the pinned Stacks Project chapter file sites-cohomology.tex; no statement or proof step was rewritten, reordered or removed.\par\medskip
\noindent\textit{Editorial note.} The proof is the author's complete argument; the reduction and the triangle are both inside the counted span.\par\medskip
\noindent\textit{Editorial note.} Every labelled reference occurring inside the counted spans is carried in assembly.external\_references as a non-load-bearing navigation pointer (same-chapter ones at their definition site, cross-chapter ones at the invocation inside the counted span), so no load-bearing mathematics is replaced by a stub. No bibliography citation occurs inside the counted spans.\par\medskip
\noindent\textit{Editorial note.} The checked chapter list for same-name collisions is the set of stacks chapters staged in this repository: chow, cohomology, divisors, etale-cohomology, flat, more-algebra, more-morphisms, morphisms, spaces-more-morphisms, varieties together with this card's chapter sites-cohomology.\par\medskip
\noindent\textit{Editorial note.} The preamble is the chapter's own retained minimal preamble (sources/stacks-sites-cohomology/retained-preamble.tex), a verbatim excerpt of the native preamble: the Xy-pic setup and the native theorem declarations with the mathematical macros.\par\medskip
\noindent\textit{Allowed standard prerequisites (editorial).}\begin{itemize}
\item Sheafification of presheaves of sets and amalgamated sums. Same chapter, the sheafification of a presheaf of sets and the description of an amalgamated sum of sheaves in terms of a pushout of presheaves. Supplies hypotheses (1) and (2) and the objects whose cohomology is compared.
\item Cohomology of a presheaf of sets as discussed in the section on lim-presheaves. Same chapter, the cohomology RGamma(K, K) attached to a presheaf of sets, defined through the topos-theoretic construction of the section. Supplies the functors appearing in the distinguished triangle.
\end{itemize}\medskip
\ifdefined\hypertarget\hypertarget{sourcepointer-1}{}\fi
\noindent\textit{Source reference.} Original section 079X: sources/stacks-sites-cohomology/pinned-original-source.tex, source lines 2017--2017. Navigation pointer invoked at proof line 6149; the counted span contains the pointer only and the cited result is not reproduced inside it, so the pointer is registered as a non-load-bearing navigation stub in the established form for Stacks material. No bibliography citation is involved.\par\medskip
\ifdefined\hypertarget\hypertarget{sourcepointer-2}{}\fi
\noindent\textit{Source reference.} Original lemma 03CI: sources/stacks-sites-cohomology/pinned-original-source.tex, source lines 6159--6159. Navigation pointer invoked at proof line 6159; the counted span contains the pointer only and the cited result is not reproduced inside it, so the pointer is registered as a non-load-bearing navigation stub in the established form for Stacks material. No bibliography citation is involved.\par\medskip
\ifdefined\hypertarget\hypertarget{sourcepointer-3}{}\fi
\noindent\textit{Source reference.} Original lemma 0EVY: sources/stacks-sites-cohomology/pinned-original-source.tex, source lines 6074--6074. Navigation pointer invoked at proof line 6167; the counted span contains the pointer only and the cited result is not reproduced inside it, so the pointer is registered as a non-load-bearing navigation stub in the established form for Stacks material. No bibliography citation is involved.\par\medskip
\begin{lemma}
\label{lemma-square-triangle-general}
Let $\mathcal{C}$ be a site. Consider a commutative diagram
$$
\xymatrix{
\mathcal{D} \ar[r] \ar[d] & \mathcal{F} \ar[d] \\
\mathcal{E} \ar[r] & \mathcal{G}
}
$$
of presheaves of sets on $\mathcal{C}$ and assume that
\begin{enumerate}
\item $\mathcal{G}^\# =
\mathcal{E}^\# \amalg_{\mathcal{D}^\#} \mathcal{F}^\#$, and
\item $\mathcal{D}^\# \to \mathcal{F}^\#$ is injective.
\end{enumerate}
Then there is a canonical distinguished triangle
$$
R\Gamma(\mathcal{G}, K) \to
R\Gamma(\mathcal{E}, K) \oplus
R\Gamma(\mathcal{F}, K) \to
R\Gamma(\mathcal{D}, K) \to
R\Gamma(\mathcal{G}, K)[1]
$$
functorial in $K \in D(\mathcal{C})$ where $R\Gamma(\mathcal{G}, -)$
is the cohomology discussed in Section \ref{section-limp}.
\end{lemma}

\begin{proof}
Since sheafification is exact and since
$R\Gamma(\mathcal{G}, -) = R\Gamma(\mathcal{G}^\#, -)$
we may assume $\mathcal{D}, \mathcal{E}, \mathcal{F}, \mathcal{G}$
are sheaves of sets. Moreover, the cohomology
$R\Gamma(\mathcal{G}, -)$ only depends on the topos,
not on the underlying site. Hence by
Sites, Lemma \ref{sites-lemma-topos-good-site}
we may replace $\mathcal{C}$ by a ``larger'' site
with a subcanonical topology such that $\mathcal{G} = h_X$,
$\mathcal{F} = h_Y$,
$\mathcal{E} = h_Z$, and
$\mathcal{D} = h_E$
for some objects $X, Y, Z, E$ of $\mathcal{C}$.
In this case the result follows from
Lemma \ref{lemma-square-triangle}.
\end{proof}

\end{document}
